\documentclass{beamer}
\usepackage{amsfonts,amsmath,oldgerm}
\usepackage{listings}
\usepackage{booktabs}
\usetheme{sintef}

\usefonttheme[onlymath]{serif}
\usetikzlibrary{positioning}      % beamer already provides tikz/pgf

\titlebackground*{assets/background}

% ---- listing style for PDDL / Prolog snippets ----
\definecolor{codecomment}{HTML}{8A7177}
\lstset{
  basicstyle=\ttfamily\small,
  commentstyle=\color{codecomment}\itshape,
  columns=fullflexible,
  keepspaces=true,
  breaklines=true,
  frame=none,
  xleftmargin=0.4em,
  morecomment=[l]{\%}            % treat % as a comment in snippets
}

% =====================================================================
%  FORMATTING FIXES (kept in main.tex so the upstream theme is unmodified)
% =====================================================================
\makeatletter
\renewcommand{\thesection}{\ifnum\value{section}>0 \arabic{section}\fi}

\renewcommand{\backmatter}[1][]{%
  \begingroup
  \themecolor{main}%
  \begin{frame}[c]
    \centering
    \begin{minipage}{\textwidth}
      \usebeamercolor[fg]{normal text}%
      \centering
      \ifstrequal{#1}{notitle}{}{%
        {\Huge\inserttitle\par}%
        \vspace{5mm}%
      }%
      {\Large\textsl{Thank you for listening! \\ Any questions?}\par}%
    \end{minipage}
  \end{frame}
  \endgroup
}

\renewenvironment{chapter}[3][]{%
  \begingroup
  \themecolor{main}%
  \setbeamertemplate{footline}{}%
  \ifstrempty{#2}{%
    \setbeamercolor{frametitle}{fg=white}%
    \setbeamercolor{normal text}{fg=white,bg=maincolor}%
  }{%
    \setbeamercolor{frametitle}{fg=white}%
    \setbeamercolor{normal text}{fg=white,bg=#2}%
  }%
  \ifstrempty{#1}{}{\setbeamertemplate{background}{\TikzSplitSlide{#1}}}%
  \setbeamertemplate{frametitle}{%
    \vspace*{8ex}%
    \begin{beamercolorbox}[wd=0.5\textwidth]{frametitle}
      \usebeamerfont{frametitle}\insertframetitle
    \end{beamercolorbox}%
  }%
  \begin{frame}{#3}
  \hspace*{0.05\textwidth}%
  \minipage{0.52\textwidth}%
  \usebeamercolor[fg]{normal text}%
  \small%
}{%
  \endminipage
  \end{frame}
  \endgroup
}
\makeatother

\title{Autonomous Library Reshelving}
\subtitle{A Planning and Reasoning Project}
\course{M.Sc. Artificial Intelligence and Robotics}
\author{Vano Mazashvili}
\IDnumber{1993251}
\date{Planning \& Reasoning --- A.A. 2025/26}

\begin{document}

\maketitle

\begin{frame}{The Domain}
\textbf{Mobile robots maintain a library by reshelving returned books.}

\vspace{0.5em}
Returned books accumulate on a central \emph{return cart}. Battery-powered
robots collect them and redistribute each to a category-matched shelf
(fiction, science, history). Navigation consumes battery; robots recharge at
\emph{charging docks}.

\vspace{0.6em}
\begin{columns}
\begin{column}{0.55\textwidth}
\begin{itemize}
  \item \textbf{Robots} --- battery-powered, mobile across zones
  \item \textbf{Books} --- each belongs to one category
  \item \textbf{Shelves} --- category-matched destinations
  \item \textbf{Docks} --- recharge points on the map
\end{itemize}
\end{column}
\begin{column}{0.42\textwidth}
\begin{block}{Goal}
Every returned book ends up shelved on a shelf of its own category --- while
the model handles battery economics and dynamic events.
\end{block}
\end{column}
\end{columns}

\vspace{0.2em}
\textbf{Two formalisms:} PDDL (classical STRIPS + numeric ENHSP) and IndiGolog
(situation calculus + controllers).
\end{frame}

\begin{frame}{Zone Topology}
The IndiGolog component uses a 7-zone graph. Edges are bidirectional; robots
navigate one hop at a time. Books start at the cart; each shelf accepts one
category; the dock is the recharge point.

\vspace{0.2em}
\begin{center}
\scalebox{0.8}{%
\begin{tikzpicture}[
    every node/.style={font=\small\bfseries, text=white},
    zone/.style={rectangle, rounded corners=3pt, minimum width=1.7cm,
                 minimum height=0.7cm, align=center},
    start/.style={zone, fill=maincolor},
    junc/.style={zone, fill=gray},
    shelf/.style={zone, fill=teal},
    dock/.style={zone, fill=green!50!black},
    edge/.style={draw=gray!55, line width=1.2pt}
  ]
  \node[start] (cart) at (0,0)   {cart};
  \node[junc]  (z1)   at (2.6,0) {z1};
  \node[shelf] (fic)  at (2.6,2) {fiction\_shelf};
  \node[junc]  (z2)   at (5.4,0) {z2};
  \node[shelf] (sci)  at (8.0,1.1){science\_shelf};
  \node[shelf] (his)  at (8.0,-1.1){history\_shelf};
  \node[dock]  (dk)   at (5.4,-2){dock};
  \draw[edge] (cart) -- (z1);
  \draw[edge] (z1)   -- (fic);
  \draw[edge] (z1)   -- (z2);
  \draw[edge] (z2)   -- (sci);
  \draw[edge] (z2)   -- (his);
  \draw[edge] (z2)   -- (dk);
\end{tikzpicture}%
}%
\end{center}

\vspace{0.1em}
{\scriptsize\color{gray}
\textbf{\color{maincolor}cart} return cart \quad
\textbf{\color{gray}gray} junction \quad
\textbf{\color{teal}teal} shelf (category) \quad
\textbf{\color{green!50!black}green} charging dock}
\end{frame}

\section{PDDL}
\begin{chapter}{maincolor}{1.\quad PDDL}
\begin{itemize}
  \item Domain: predicates \& actions
  \item A key design decision (navigation \& battery)
  \item Three problem instances
  \item STRIPS heuristic comparison \& findings
  \item Numeric (ENHSP) formulation \& findings
\end{itemize}
\end{chapter}

\begin{frame}[fragile]{STRIPS Domain --- Predicates \& Actions}
\begin{columns}
\begin{column}{0.55\textwidth}
\textbf{Fluents (change over time)}
\begin{lstlisting}
robot-at(r,z)    robot r in zone z
book-at(b,z)     book b in zone z
carrying(r,b)    r carries book b
shelved(b)       b correctly shelved
battery-low(r)   r's battery is low
free(r)          r carries nothing
\end{lstlisting}
\vspace{0.3em}
\textbf{Rigid (static)}
\begin{lstlisting}
path(z1,z2)          zones adjacent
book-category(b,c)   b has category c
shelf-category(z,c)  z accepts c
station-at(z)        z has a dock
\end{lstlisting}
\end{column}
\begin{column}{0.42\textwidth}
\begin{block}{Actions}
\begin{description}
  \item[navigate(r,z1,z2)] drive along a path edge
  \item[pickup(r,b)] collect a co-located book
  \item[shelve(r,b)] place on matching shelf --- \emph{depletes battery}
  \item[recharge(r)] restore battery at a dock
\end{description}
\end{block}
\end{column}
\end{columns}
\end{frame}

\begin{frame}{Design Decision --- Navigation \& Battery}
\small
\textbf{\large \alert{In the STRIPS and IndiGolog models, navigate has no battery precondition.}}

\vspace{0.5em}
Requiring a charged battery to move would create \textbf{unrecoverable
dead-ends}: a robot could strand itself at a shelf with a low battery and no
way to reach a dock. So navigation is always permitted along an existing edge.
The numeric model is the exception --- it \emph{can} safely guard navigate
(battery $\geq 10$) because its planner reasons about quantities and plans
recharges ahead; the boolean models cannot, so they omit the precondition.

\vspace{0.3em}
\begin{columns}
\begin{column}{0.48\textwidth}
\begin{block}{Why it matters}
Battery as an action precondition can make a goal physically unreachable ---
the planner reports failure not because no plan exists in principle, but
because the model painted the robot into a corner.
\end{block}
\end{column}
\begin{column}{0.48\textwidth}
\begin{colorblock}[white]{green!50!black}{How each model handles it}
\textbf{STRIPS} --- battery is an abstract boolean flag; \texttt{shelve}
depletes it, \texttt{recharge} clears it.\\[0.4em]
\textbf{Numeric (ENHSP)} --- battery is a quantity (0--100); the planner
balances delivery against recharging trips.
\end{colorblock}
\end{column}
\end{columns}
\end{frame}

\begin{frame}{Three Problem Instances}
Each instance increases in complexity, exercising the planners at different
scales.

\vspace{0.8em}
\begin{center}
\begin{tabular}{lccccc}
\toprule
 & \textbf{Robots} & \textbf{Zones} & \textbf{Books} & \textbf{Categories} & \textbf{Docks}\\
\midrule
\textbf{P1} & 1 & 5  & 2 & 2 & 1\\
\textbf{P2} & 2 & 8  & 4 & 3 & 1\\
\textbf{P3} & 2 & 12 & 6 & 4 & 2\\
\bottomrule
\end{tabular}
\end{center}

\vspace{0.8em}
{\small\color{gray}\itshape Topology is a connected grid; P3 adds a second
dock and a fourth category, pushing the optimal-search configurations to their
limits.}
\end{frame}

\begin{frame}{STRIPS --- Heuristic Comparison (A$^\star$)}
Fast-Downward, A$^\star$ with four heuristics. Plan cost $=$ length (unit
actions). Expansions shown per problem.

\vspace{0.6em}
\begin{center}
\small
\begin{tabular}{lcccccc}
\toprule
\textbf{Heuristic} & \textbf{P1 len} & \textbf{P1 exp} & \textbf{P2 len} & \textbf{P2 exp} & \textbf{P3 len} & \textbf{P3 exp}\\
\midrule
Blind          & 13 & 45 & 27 & 17{,}355 & 44 & 473{,}273\\
$h^{\max}$     & 13 & 37 & 27 & 15{,}082 & 44 & 447{,}251\\
$h^{add}$      & 13 & 16 & 27 & 3{,}893  & \textcolor{maincolor}{\textbf{46$^\star$}} & 106{,}880\\
$h^{FF}$       & 13 & 26 & 27 & 9{,}496  & 44 & 346{,}822\\
\bottomrule
\end{tabular}
\end{center}

\vspace{0.5em}
{\small\itshape\color{maincolor}$^\star$ $h^{add}$ returns a 46-step plan on
P3 --- suboptimal (optimum is 44). Detailed next slide.}
\end{frame}

\begin{frame}{STRIPS --- Key Findings (Problem 3)}
\small
\begin{block}{Inadmissibility, observed}
$h^{add}$ returned a 46-step plan where Blind, $h^{\max}$ and $h^{FF}$ all
found the 44-step optimum. By overestimating, A$^\star$ committed early to a
non-optimal path. ($h^{FF}$ is also inadmissible but stayed optimal here ---
luck, not guarantee.)
\end{block}
\vspace{-0.5em}
\begin{block}{Informativeness $\neq$ speed}
$h^{\max}$ and $h^{FF}$ were \emph{slower} than Blind in wall-clock (2.30\,s,
2.36\,s vs 1.19\,s) despite fewer expansions: on this grid topology,
delete-relaxation heuristics are loose, so per-node cost outweighs pruning.
\end{block}
\begin{block}{Cost-model artifact}
Under unit action costs, dock\_2 is never used in P3 --- dock\_1's central
position dominates. The numeric model (differential costs) is where dock\_2
could earn its place.
\end{block}
\end{frame}

\begin{frame}[fragile]{Numeric PDDL --- Battery as a Fluent}
Per Prof.\ Marrella's recommendation, an additional numeric version models
battery quantitatively with ENHSP.

\vspace{0.6em}
\begin{columns}
\begin{column}{0.46\textwidth}
\begin{block}{Model}
\begin{lstlisting}
battery    numeric 0-100
navigate   -10 batt, +1 cost
shelve     -30 batt, +2 cost
recharge   set 100, +3 cost
metric     minimize total-cost
\end{lstlisting}
\end{block}
\end{column}
\begin{column}{0.50\textwidth}
\begin{colorblock}[white]{green!50!black}{What it captures}
Unlike the STRIPS flag, the numeric model expresses real battery economics:
navigation costs energy, and the planner must weigh delivery against
recharging trips. This answers a different question --- not the shortest plan,
but the most energy-efficient one.
\end{colorblock}
\end{column}
\end{columns}
\end{frame}

\begin{frame}{ENHSP --- Configuration Comparison}
Four configurations across the three problems. Satisficing (\texttt{sat-}) vs.\
optimal (\texttt{opt-}) search.

\vspace{0.6em}
\begin{center}
\small
\begin{tabular}{lcccc}
\toprule
\textbf{Config} & \textbf{P1 cost} & \textbf{P2 cost} & \textbf{P3 cost} & \textbf{P3 time}\\
\midrule
\texttt{aibr}      & 20 (subopt) & ---  & ---     & ---\\
\texttt{sat-hadd}  & 15 & 34 & 64 & 0.27\,s\\
\texttt{opt-hmax}  & 15 & 31 & \textcolor{maincolor}{\textbf{timeout}} & \textcolor{maincolor}{\textbf{$>$120\,s$^\star$}}\\
\texttt{opt-hrmax} & 15 & 31 & \textcolor{maincolor}{\textbf{timeout}} & \textcolor{maincolor}{\textbf{$>$120\,s$^\star$}}\\
\bottomrule
\end{tabular}
\end{center}

\vspace{0.5em}
{\small\itshape\color{maincolor}$^\star$ Optimal configs do not terminate
within 120\,s on P3, while satisficing solves in 0.27\,s.}
\end{frame}

\begin{frame}{Numeric --- Key Finding}
\begin{block}{Exponential blow-up in optimal numeric planning}
On Problem 3, \texttt{sat-hadd} and \texttt{sat-hff} return plans in 0.27\,s,
while \texttt{opt-hmax} and \texttt{opt-hrmax} fail to terminate within
120\,s. This is the empirical signature of the cost of admissibility:
guaranteeing optimality over a numeric state space is dramatically more
expensive than finding a good-enough plan.
\end{block}

\vspace{0.3em}
\begin{block}{STRIPS vs. Numeric --- two different questions}
\textbf{STRIPS} asks ``what is the shortest plan?'' --- P3 optimum 44 steps in
$\sim$1.2\,s.\\[0.4em]
\textbf{Numeric} asks ``what is the most energy-efficient plan?'' --- P3
satisficing 52 steps in 0.27\,s, with explicit battery accounting the STRIPS
model cannot express.
\end{block}
\end{frame}

\section{IndiGolog}
\begin{chapter}{maincolor}{2.\quad IndiGolog}
\begin{itemize}
  \item Situation-calculus formulation (the BAT)
  \item Fluents, actions, causal laws
  \item Reasoning tasks: legality \& projection
  \item Controller 1 --- basic search
  \item Controller 2 --- reactive (prioritized interrupts)
\end{itemize}
\end{chapter}

\begin{frame}{Situation-Calculus Formulation}
\footnotesize
The same domain is recast as a \textbf{Basic Action Theory} and implemented on
the standalone \texttt{indigolog\_plain} interpreter.

\vspace{-0.4em}
\begin{block}{Interpreter}
\texttt{indigolog\_plain} --- the vanilla architecture used by the course's
elevator examples. It natively supports exogenous actions, \texttt{search},
and \texttt{prioritized\_interrupts}; no deployment engine, devices, or
networking needed.
\end{block}
\vspace{-0.6em}
\begin{block}{Single robot}
The PDDL component already demonstrates multi-robot coordination (P2, P3).
IndiGolog focuses on the reasoning tasks and control constructs, so a single
agent keeps controller traces legible --- a deliberate, justified adaptation.
\end{block}
\vspace{-0.6em}
\begin{block}{Boolean battery}
Battery is a boolean here, set only by an exogenous drain event and cleared by
\texttt{recharge}. The numeric treatment lives in the PDDL part, exactly where
the recommendation scoped it.
\end{block}
\end{frame}

\begin{frame}[fragile]{BAT --- Fluents, Actions, Causal Laws}
Written in \texttt{indigolog\_plain} style: \texttt{prim\_fluent},
value-equality conditions, \texttt{causes\_val} successor-state axioms.

\vspace{0.1em}
\begin{lstlisting}[basicstyle=\ttfamily\scriptsize]
% fluents (boolean modelled as value-equality, never bare atoms)
prim_fluent(robot_loc(R)) :- robot(R).
prim_fluent(battery_low(R)) :- robot(R).

% precondition: shelve needs a charged battery and a matching shelf
poss(shelve(R,B), and(carrying(R)=B,
                  and(robot_loc(R)=Z,
                  and(shelf_category(Z,C),
                  and(book_category(B,C), battery_low(R)=false))))).

% causal laws - navigate does NOT deplete battery (see design decision)
causes_val(navigate(R,_,To), robot_loc(R), To, true).
causes_val(battery_drained(R), battery_low(R), true, true).
causes_val(recharge(R), battery_low(R), false, true).
\end{lstlisting}
\end{frame}

\begin{frame}[fragile]{Reasoning Tasks --- Legality \& Projection}
\begin{columns}
\begin{column}{0.49\textwidth}
\begin{block}{1 --- Legality}
{\footnotesize Is a sequence executable from $S_0$?}
\begin{lstlisting}
legal_seq([
  navigate(r1,cart,z1),
  navigate(r1,z1,fiction_shelf)]).
  -> true

legal_seq([
  navigate(r1,cart,z1),
  pickup(r1,b1)]).
  -> false
\end{lstlisting}
\end{block}
\end{column}
\begin{column}{0.49\textwidth}
\begin{block}{2 --- Projection}
{\footnotesize Does a fluent hold after a sequence?}
\begin{lstlisting}
holds_after(
  [navigate(r1,cart,z1)],
  robot_loc(r1)=z1).
  -> true

holds_after(
  [navigate(r1,cart,z1)],
  battery_low(r1)=false).
  -> true
\end{lstlisting}
\end{block}
\end{column}
\end{columns}
\end{frame}

\begin{frame}[fragile]{Controller 1 --- Basic Search}
A goal-directed, bounded program wrapped in the IndiGolog \texttt{search}
operator finds a complete plan offline. A visited-set in routing forces simple
paths.

\vspace{0.1em}
\begin{columns}
\begin{column}{0.48\textwidth}
\begin{block}{\texttt{main(basic).} $\to$ 19 actions}
\begin{lstlisting}[basicstyle=\ttfamily\footnotesize]
pickup(r1,b1)
navigate(r1,cart,z1)
navigate(r1,z1,fiction_shelf)
shelve(r1,b1)
navigate(r1,fiction_shelf,z1)
navigate(r1,z1,cart)
pickup(r1,b2)
...
shelve(r1,b3)
\end{lstlisting}
\end{block}
\end{column}
\begin{column}{0.48\textwidth}
\textbf{Constructs:} \texttt{search} (offline lookahead), \texttt{pi}/\texttt{ndet}
(nondeterministic choice), \texttt{while} (loop until all shelved).
\begin{colorblock}[white]{maincolor}{Finding: a plan, not the best}
The first working version wandered (35 actions). A visited-set forbidding
revisits forces acyclic routes $\to$ 19 actions. Same lesson as the PDDL
$h^{add}$ finding: quality needs structure or an admissible heuristic.
\end{colorblock}
\end{column}
\end{columns}
\end{frame}

\begin{frame}[fragile]{Controller 2 --- Reactive (Prioritized Interrupts)}
\texttt{prioritized\_interrupts} react to three exogenous events. Priority
arbitrates between them; each produces a genuine behavioural response.

\vspace{0.1em}
\begin{columns}
\begin{column}{0.46\textwidth}
\begin{description}
  \item[battery\_drained(r)] divert to dock $\to$ recharge $\to$ resume
  \item[urgent\_request(b)] serve book $b$ first (pre-empt)
  \item[new\_return(b)] book reappears at cart $\to$ reshelve
\end{description}
\end{column}
\begin{column}{0.50\textwidth}
\begin{block}{\texttt{battery\_drained(r1)} injected}
\begin{lstlisting}[basicstyle=\ttfamily\footnotesize]
start_interrupts
navigate(r1,cart,z1)
navigate(r1,z1,z2)
navigate(r1,z2,dock)
recharge(r1)        <- reaction
navigate(r1,dock,z2)
...
shelve(r1,b3)
stop_interrupts
\end{lstlisting}
\end{block}
\end{column}
\end{columns}
\end{frame}

\begin{frame}{Cross-Component Synthesis}
\small
\textbf{One domain, three complementary treatments}
\vspace{0.1em}
\begin{columns}
\begin{column}{0.32\textwidth}
\begin{colorblock}[white]{maincolor}{STRIPS}
battery $=$ boolean flag\\ {\small planning over an abstraction}
\end{colorblock}
\end{column}
\begin{column}{0.32\textwidth}
\begin{colorblock}[white]{teal}{Numeric}
battery $=$ quantity 0--100\\ {\small energy economics, ENHSP}
\end{colorblock}
\end{column}
\begin{column}{0.32\textwidth}
\begin{colorblock}[white]{green!50!black}{IndiGolog}
battery $=$ reactive event\\ {\small drive-to-dock \& recharge}
\end{colorblock}
\end{column}
\end{columns}

\vspace{0.15em}
{\footnotesize\textbf{\color{maincolor}Depletion differs by design:}
\textit{STRIPS auto-depletes on \texttt{shelve}; IndiGolog depletes only via
the exogenous \texttt{battery\_drained} event --- the reactive controller needs
an external trigger.}}

\vspace{0.1em}
\begin{columns}
\begin{column}{0.48\textwidth}
\begin{block}{Suboptimality, twice}
PDDL: $h^{add}$ returns 46 vs the 44 optimum (inadmissibility). IndiGolog:
naive search wanders 35 steps until a visited-set forces 19. Same lesson ---
quality needs an admissible heuristic or a structural constraint.
\end{block}
\end{column}
\begin{column}{0.48\textwidth}
\begin{block}{Reactive vs. planned}
The search controller is globally optimal but offline. The reactive controller
is responsive but locally greedy --- it recharges, then returns to the cart, a
visible non-optimal detour. Two paradigms, two trade-offs.
\end{block}
\end{column}
\end{columns}
\end{frame}

\begin{frame}{Adaptations Justified at Presentation}
Per the project rules, the domain and problem instances stay fixed; component
divergences are declared here:
\begin{itemize}
  \item \textbf{Single robot} in IndiGolog vs.\ up to two in PDDL --- focus on
        reasoning tasks and control constructs, not coordination.
  \item \textbf{Boolean battery} in IndiGolog vs.\ numeric in PDDL --- numeric
        treatment scoped to the PDDL part, as recommended.
  \item \textbf{Exogenous events} (\texttt{urgent\_request}, \texttt{new\_return},
        \texttt{battery\_drained}) have no PDDL counterpart --- they exercise
        \texttt{prioritized\_interrupts}.
\end{itemize}

\vspace{0.6em}
\begin{center}
{\small Repository: \texttt{github.com/vmazashvili/planning-reasoning-library-reshelving}}
\end{center}
\end{frame}

\backmatter

\end{document}